\begin{textsample}{POS Dim 4 – pa – Score 15.00 – pa/pa\_em\_53.txt}  \label{ex:f4_pos_006}
O poder, a jurisdição, a administração e a soberania obedecem aos princípios de \textbf{fronteira}.
Ou seja, historicamente, ao expressar as diversas culturas, grupos definem \textbf{fronteiras}, formas de organização social e, por vezes, áreas de confronto com outros grupos.
O que também leva a ampliar o conceito de \textbf{fronteira}, não somente entendendo -o como os limites espaciais entre \textbf{territórios}, mas também como \textbf{fronteiras} culturais e econômicas que igualmente definem diferenças, \textbf{exclusões} e \textbf{identidades} perante os grupos sociais ou econômicos.
Tais processos podem ser percebidos, por exemplo, na formação dos \textbf{estados} nacionais, bem como nas lutas pela posse da terra, ou mesmo através da cultura dos povos amazônidas, os quais expressam a sua \textbf{identidade}, suas \textbf{territorialidades} e sua \textbf{fronteira} a partir de elementos intangíveis, mas que dão forma à sua maneira de habitar no mundo.
Ao contextualizar essas categorias com as realidades regionais, Petit ( 2003 ) considera a categoria \textbf{território} e \textbf{fronteira} como essencial para entender as problemáticas que envolvem o processo de colonização e ocupação da Amazônia.
Existe uma importante relação entre \textbf{território}, economia e política, mostrando como localmente se deram as problemáticas de uma região pensada como “ área de \textbf{fronteira} ” até meados dos anos 1970 ( PETIT, 2003 ; BECKER, 2004 ).
Ainda segundo o autor, é necessário refletir sobre a questão do \textbf{território} e da \textbf{fronteira} para compreender a configuração espacial da região ao longo do tempo.
Petit ( 2003, p.32 ) afirma que o \textbf{território} deve ser entendido como “ vinculado à apropriação, controle ou domínio exercido numa determinada área ” e o conceito também está relacionado com “ os recortes do \textbf{território} que fixam as \textbf{fronteiras} entre países e, também, os limites regionais e divisões político-administrativas internas dos Estados-Nação”.[^90 ] Assim, a ampliação das \textbf{fronteiras} nacionais para a Amazônia trouxe uma série de conflitos e novas dinâmicas espaciais e, por essa razão, precisamos pensar o papel de inúmeros indivíduos nessas sociabilidades, bem como pensar o lugar dos excluídos dos processos de expansão e de colonização realizados ao longo dos séculos por grupos que se deslocaram para a região e incorporaram esses \textbf{territórios} ( LOUREIRO, 1992 ).
Outra questão que se remete à análise a partir desta categoria \textbf{território} e \textbf{fronteira} é a violência no campo em seu contexto regional, levando em consideração que a presença de novas empresas e \textbf{capitais} alteraram em grande medida a vida das diferentes comunidades que habitam o \textbf{Estado} do Pará.
É importante entender o ponto de vista de diferentes sujeitos envolvidos com a luta no campo, em especial, como os processos de colonização ocorridos em diferentes momentos causaram inúmeros conflitos agrários, fundamentados na disputa pelo \textbf{território}, pensado como um elemento central na produção econômica do interior do Pará ( SACRAMENTO, 2012 ).
A BNCC estabeleceu a categoria indivíduo, natureza, sociedade, cultura e ética como mais uma balizadora da área CHSA ; a mesma resultou na junção de dimensões muito vastas enquanto base da referida área ( BRASIL, 2018 ).
Contudo, neste documento curricular, com o intuito de reconhecer e valorizar as peculiaridades do ensino médio do Pará, além de lançar um olhar crítico e propositivo para os desdobramentos conceituais de cada elemento integrador das CHSA, surge, então, a necessidade de desmembrar a categoria supracitada.
Propõe -se, portanto, o seguinte redimensionamento : Natureza e Cultura, como primeiro binômio ; e Sociedade, Indivíduo, \textbf{Identidade} e Interculturalidade, enquanto segundo desdobramento, considerando sua intrínseca relação.
Entende -se que essa organização estrutural visa facilitar o trabalho docente ao desmembrar a grande categoria originalmente presente na BNCC.
Por outro lado, é necessário frisar que não se defende que sejam edificadas dicotomias entre as categorias criadas, ao contrário, elas devem ser constantemente problematizadas na sua abordagem em sala de aula.
Portanto, a terceira categoria de área é Natureza e Cultura, fazendo referência à grande dualidade que compõe a própria existência humana no mundo.
Em primeiro lugar, revelando que a natureza é um tema que, por si só, já se configura como uma dimensão de extrema relevância para qualquer problematização que se efetue a respeito das relações que o ser humano estabelece no mundo, pois se destaca enquanto força geradora.
Ela, por um lado, se impõe à revelia dos ditames humanos, e, por outro, padece com a exploração desencadeada pelos mesmos agentes que a transformam em recurso disponível para o seu usufruto.
Essa diversificação da natureza é processo e resultado das relações estabelecidas no tempo e no espaço ( SANTOS, 2008 ).
Assim, cabe destacar a relação do \textbf{homem} ( eu ) com os outros e enquanto produtor de cultura, bem como entender a referida dinâmica de apropriação da natureza como elemento de produção e reprodução material e simbólica da sociedade em diferentes escalas histórico-geográficas.
Além disso, é essencial ressaltar as especificidades do \textbf{território} e dos povos \textbf{amazônicos}, o que nos permite refletir sobre as singularidades humanas em suas várias dimensões, propondo uma visão crítica de certos elementos que naturalizam desigualdades históricas.
Essas categorias são fundamentais, por exemplo, para problematizar representações \textbf{amazônicas} ao situá -las historicamente.
Podemos perceber, entre outros fenômenos, uma mudança a partir de meados do século XIX na maneira de conceber a “ Amazônia ”, não apenas sob a ótica do mundo natural, mas também sob um olhar mais sociológico, filosófico, histórico, geográfico e \textbf{antropológico} ( FIGUEIREDO, 2017 ).
Desse modo, ao refletir no espaço escolar sobre tais conceitos e a partir dessa categoria, pode -se contribuir à compreensão das diferentes representações e noções sobre a Amazônia enquanto um constructo histórico.
Ou seja, preparar caminhos para compreender a complexa relação entre cultura e natureza, em que se possa desenvolver a respectiva crítica às tradições que persistem em compreender a região \textbf{amazônica} a partir de uma visão edênica e sem a complexa relação histórica estabelecida entre múltiplos sujeitos ao longo do tempo ( FREIRE, 2019 ).
Como já exposto na concepção deste Documento, devemos reiterar que há uma diversidade significativa no âmbito das concepções de Amazônia.
Por essa razão, parte -se da premissa de que existe uma diversidade de Amazônias[^91 ], o que implica em problematizar a região a partir das suas diferenças, suas particularidades, suas \textbf{territorialidades} em uma perspectiva intercultural ( GONÇALVES, 2012 ).
Neste sentido, na região \textbf{amazônica} há uma relação homens/mulheres-natureza muito específica, na medida em que desenvolvem modos de vida únicos, podendo inclusive, coabitar com culturas completamente diferentes entre si, os quais acarretam na necessidade de pensar a referida categoria de área, relacionando -a com a heterogeneidade dos povos, comunidades, \textbf{identidades} e práticas cotidianas envolvidas com o seu ambiente natural nas Amazônias paraenses.
Desse modo, entendendo os povos da \textbf{floresta} como agentes históricos-sociais e portadores de saberes particulares elaborados de maneira ancestral e tradicional, assim, configura -se o \textbf{estado} do Pará, com uma pluralidade única de conhecimentos, oriundos de múltiplas matrizes, as quais precisam ser conhecidas e respeitadas em sua integralidade.
Assim, percebe -se a fecunda relação desta categoria com os três princípios curriculares da educação básica paraense e a partir destes com outras áreas do conhecimento como : linguagens e suas tecnologias e ciências da natureza e suas tecnologias, configurando -se importante elemento integrador e interdisciplinar. [ ^90 ] : Desde sua ocupação territorial pós década de 1950, o espaço \textbf{amazônico} vem sendo palco da materialização de várias \textbf{territorialidades} ligadas a lógica global de apropriação dos recursos naturais e do conhecimento dos povos da \textbf{floresta}.
Em alguns casos tal fato se deu pelas concessões governamentais, as quais foram agentes ( passivos ) deste processo/modelo de apropriação da biodiversidade e da expropriação de \textbf{homens} e \textbf{mulheres} \textbf{amazônicos}.
Os agentes de gestão pública deveriam elaborar e executar políticas públicas mais abrangentes que pudessem amenizar as principais problemáticas dos povos da \textbf{floresta}, sejam problemáticas políticas, econômicas, sociais, culturais, ambientais e tecnológicas. para Bertha Becker as intervenções governamentais tem que estar conectadas com as \textbf{territorialidades} dos indivíduos \textbf{amazônicos} e levar em “ dois vetores de transformação regional, que expressam a estrutura transicional do \textbf{Estado} e do \textbf{território} contemporâneos, o vetor tecnoindustrial e o vetor tecno-ecológico ” ( BECKER, 2010 ).
Agindo assim os governos serão mais coerentes com as demandas do \textbf{território} \textbf{amazônico}. [ ^91 ] : Termo criado pelo geógrafo Carlos Walter Porto Gonçalves, com o intuito de problematizar sobre a existência de várias Amazônias dentro do \textbf{território} \textbf{amazônico}.
É, acima de tudo, uma análise geográfica que discute as diversidades naturais, sociais, econômicas, culturais/antropológicas/territoriais e políticas da região.
O núcleo duro deste debate é afirmar a (r)existência dos atores territoriais, suas coletividades e \textbf{territorialidades} \textbf{amazônicas} ( \textbf{índios}, quilombolas, seringueiros, \textbf{ribeirinhos}, parteiras, garimpeiros, castanheiros, caboclos, \textbf{imigrantes} e outros ).
GONÇALVES, Carlos Walter Porto.
Amazônia, Amazônias.
São Paulo : Contexto, 2005.

% matched lemmas: amazônico, antropológico, capital, estado, exclusão, floresta, fronteira, homem, identidade, imigrante, mulher, ribeirinho, territorialidade, território, índio
\end{textsample}
